\input{./._relpath-to-latexroot.ltx}
\documentclass[\latexroot/\projectname]{subfiles}
\input{\latexroot/@resources/subfile-setup.ltx}
\begin{document}

\webonly{\newcommand{\versn}{Web}}

\title{Entrepreneurship, Frictions, and Wealth}

\author{Rudan Zheng}

\keywords{entrepreneurship, wealth inequality, borrowing constraints, occupational choice, bequests, reproduction}

\jelclass{D31, E21, J23 \\[0pt]
  \href{https://econ-ark.org}{\includegraphics[width=1.5cm]{images/PoweredByEconARK}}
}

\date{\today}

\maketitle
\vspace{-0.5\baselineskip}

\whenintegrated{\label{abstract}}
\begin{verbatimwrite}{Cagetti2005tk-Abstract.txt}
  Cagetti and De Nardi (2006) study how entrepreneurship, borrowing constraints, and bequests contribute to wealth concentration in the United States. We solve their household problem at fixed prices. Households choose between wage work and entrepreneurship, and entrepreneurs choose business capital subject to an enforcement constraint. We use the endogenous grid method for consumption and saving. We use a one-sided derivative at the upper asset-grid endpoint. The baseline value iteration meets its stopping rule after 115 iterations. In our young-only distribution, 8.33 percent of households choose entrepreneurship and the wealth Gini is 0.67. The original paper reports 7.5 percent and 0.80 in its full stationary equilibrium. Our distribution excludes aging and inheritance flows, and we take prices and calibrated parameters as given.
\end{verbatimwrite}
\begin{abstract}
  \input{Cagetti2005tk-Abstract.txt}
\end{abstract}
\vspace{-0.5\baselineskip}

\whenintegrated{\label{links}}
\begin{footnotesize}
  \setlength{\parskip}{0pt}
  \parbox{0.9\textwidth}{
    \begin{tabbing}
      \texttt{~~~~~~~~~~~~~~~~~~~~~~~} \= \=  \\
      \texttt{~~~~GitHub~Repo:~~~~~~~} \> \> \texttt{\href{https://github.com/\owner/Cagetti2005tk}{https://github.com/\owner/Cagetti2005tk}} \\
      \texttt{~~~~~~~~~~~~~~~~~~PDF:~} \> \> \texttt{\href{https://github.com/\owner/Cagetti2005tk/blob/main/Cagetti2005tk.pdf}{Cagetti2005tk.pdf}} \\
    \end{tabbing}
  }
\end{footnotesize}
\vspace{-0.3\baselineskip}

\centerline{Original article: Cagetti and De Nardi (2006), Journal of Political Economy. \href{https://doi.org/10.1086/508032}{DOI: 10.1086/508032}}
\vspace{-0.3\baselineskip}

\begin{authorsinfo}
  \noindent \name{This is a student reproduction analysis paper based on the original article by Marco Cagetti and Mariacristina De Nardi.}
\end{authorsinfo}

\pagenumbering{gobble}

\whenintegrated{\label{funding-info}}
\vspace{0.3\baselineskip}
\begin{minipage}{\textwidth}
  \footnotesize This project follows the course reproduction format and focuses on the theoretical household problem from Cagetti and De Nardi (2006).
  \input{\latexroot/@local/acknowledgments.texinput}
\end{minipage}

\titlepagefinish
\smartbib

\end{document}
